Este trabalho apresenta o EDGAR, uma arquitetura de Geração Aumentada por Recuperação Multimodal (MRAG) de granularidade fina, voltada à consulta de documentos didáticos em PDF. A arquitetura combina técnicas de análise de documentos, extração textual e descrição de elementos visuais para organizar e recuperar informações em nível de componentes. Essa abordagem permite utilizar conteúdos textuais e visuais na construção do contexto das respostas, preservando a rastreabilidade das evidências no documento original. Para demonstrar seu funcionamento, foi implementada uma interface web interativa que permite a consulta aos materiais e a visualização das evidências recuperadas. A arquitetura foi submetida a uma avaliação experimental comparativa com abordagens tradicionais de RAG textual, considerando a recuperação de informações e a qualidade das respostas geradas. Os resultados evidenciam contribuições do enriquecimento visual nas consultas que dependem de figuras e de raciocínio multimodal, bem como na pontuação média de qualidade das respostas, embora não indiquem superioridade geral sobre a indexação em fragmentos textuais na localização inicial de páginas. Nesse sentido, o EDGAR visa apoiar interações mais fundamentadas e rastreáveis entre o estudante e o material didático, facilitando o acesso às informações textuais e visuais presentes nos documentos.

\palavraschave{geração aumentada por recuperação; recuperação multimodal; granularidade fina; processamento de documentos; rastreabilidade visual.}
